% !TEX root = 06-rvole.tex
% title: 随机 VOLE


\section{铺垫}

回顾 SoftSpokenOT，其 Sender 得到 $L$ 对 OT 密钥，Receiver 得到相应的 $L$ 个单侧密钥。
对第 $j$ 实例，$j \in[0, L)$：

\begin{itemize}
\item
  Sender 持有成对密钥 $\rho^0_j, \rho^1_j$，不知 Receiver 选哪一边。
\item
  Receiver 持有选择位 $\beta_j\in\left\{0,1\right\}$，
  所选密钥 $\rho^{\beta_j}_j$，不知另一个密钥 $\rho_j^{1-\beta_j}$。
\end{itemize}

记号约定：$\vec\beta \in \left\{0,1\right\}^L$ 指选项串本身（即 SoftSpokenOT 的 $\beta$），
而不带箭头的 $\beta := \sum_{j\in[0,L)}\beta_j g_j \pmod n$ 专指它的 gadget 加权聚合。

拿到这些密钥后，兑现随机 MtA 关系 $y + z := w\beta \pmod n$。其中 Sender 输入 $w$，输出 $y$；
Receiver 输入 $\beta$，输出 $z$。注意本协议不兑现 ECDSA MtA 关系 $y + z := wx \pmod n$。

本协议每次使用新鲜随机的 $\vec\beta$，再由 ECDSA 编排层转换为所需的乘法关系。
安全性尚待理解，暂不补写安全论证。
协议原文见\href{https://eprint.iacr.org/2023/765.pdf}{DKLS23 Protocol 5.2}。

兑现 MtA 关系有以下两种思路：

\subsection{朴素路线：对消息进行加密}

Sender 构造 OT 消息：
\begin{align}M^0_j &:= r_j \stackrel{\$}{\leftarrow}\mathbb{Z}_n,\\
M^1_j &:= r_j + w\cdot g_j \pmod{n}.
\end{align}

把密钥直接加到消息上得到密文：
\begin{align}C^0_j &:= \rho^0_j + M^0_j, \\
C^1_j &:= \rho^1_j + M^1_j.
\end{align}

Receiver 用所持的 $\rho^{\beta_j}_j$ 解开 $C^{\beta_j}_j$，得到所选的消息：
\begin{equation}\label{eq:06-rvole:1}
M^{\beta_j}_j = r_j + \beta_j\cdot w\cdot g_j.
\end{equation}

Sender 和 Receiver 分别聚合出自己的加法秘密份额。聚合采用 gadget 方式，
详见 \hyperref[sec:gadget-vector]{「杂记：gadget 向量」}：
\begin{align}y &:= -\sum_j  r_j \pmod{n}, \\
z &:= \sum_j M^{\beta_j}_j \pmod{n}.
\end{align}

协议至此结束。通信量：对每个 OT 槽位，Sender 发出 2 个密文 $C^0_j, C^1_j$。

\subsection{另类路线：把密钥直接解读为随机数}

既然消息和密钥都是随机的，
我们把 $\rho_j^0, \rho_j^1$ 扩展为 $\mathbb{Z}_n$ 标量，如公式 \eqref{eq:06-rvole:alphas}。
这里的扩展及域采样遵循\hyperref[sec:rvole:parameters-model]{完全版的参数与模型约定}，不能把短密钥直接取模当作均匀域元素。
\begin{equation}\label{eq:06-rvole:alphas}
\alpha_j^b := \mathrm{PRF}^{\mathbb Z_n}_{\rho_j^b}\bigl(\mathrm{Enc}(\mathtt{sid},\texttt{"ot-pad"},j,0)\bigr),\quad b\in\{0,1\}.
\end{equation}

Sender 拥有 $\alpha_j^0,\alpha_j^1$，以及秘密因子 $w$。
Receiver 拥有 $\alpha_j^{\beta_j}$，以及秘密因子 $\beta$。

对每个 OT 实例 $j$，也就是 gadget 分解后的第 $j$ 分量，Sender 发送修正量：
\begin{equation}\label{eq:06-rvole:3}
\tilde a_j := \alpha^0_j - \alpha^1_j + w \pmod{n}.
\end{equation}

最后，Sender 按公式 \eqref{eq:06-rvole:y} 得到加法秘密份额 $y$。
类似地，Receiver 按公式 \eqref{eq:06-rvole:z} 得到加法秘密份额 $z$。
注意到 $\alpha^{\beta_j}_j + \beta_j\tilde a_j = \alpha^0_j + \beta_j w \pmod n$，
读者凭此公式自行验算公式 \eqref{eq:06-rvole:y} 和 \eqref{eq:06-rvole:z}。
\begin{align}\label{eq:06-rvole:y}
y &:= -\sum_j g_j\cdot\alpha^0_j \pmod{n}, \\
\label{eq:06-rvole:z}
z &:= \sum_j g_j\left(\alpha^{\beta_j}_j + \beta_j\tilde a_j\right) \pmod{n}.
\end{align}

协议至此结束。通信量：对每个 OT 槽位，Sender 发送 1 个标量 $\tilde a_j$。
修正量 $\tilde a_j$ 的构造不依赖 gadget 权重 $g_j$；最终份额的聚合仍需使用 $g_j$。

\begin{center}\rule{0.5\linewidth}{0.5pt}\end{center}

\section{RVOLE 协议核心版}

RVOLE 完全版有三个索引维度：OT 实例 $j$，检查编号 $k$，MtA
关系编号 $i$ 或 $k$。维度太多会干扰认知。本章描述单个负载且无检查的
RVOLE 协议。

\subsection{规格}

【输入】Sender 持有秘密 $w$，
Receiver 持有秘密 $\beta=\sum_{j\in[0,L)} \beta_j \cdot g_j$。

【输出】Sender 获得 $y$，Receiver 获得 $z$，使得
\begin{equation}\label{eq:06-rvole:4}
y + z := w \cdot \beta \pmod{n}.
\end{equation}

\subsection{实施}
\subsubsection{Sender 发送修正项，获得加法份额}

Sender 把 OT 密钥 $\rho^0_j$ 转换为随机标量 $\alpha^0_j\in\mathbb{Z}_n$。
同理，把 $\rho^1_j$ 转换为随机标量 $\alpha^1_j\in\mathbb{Z}_n$。
然后构造修正项：
\begin{equation}\label{eq:06-rvole:5}
\tilde{a}_{j} = \alpha^0_{j} - \alpha^1_{j} + w \pmod{n}.
\end{equation}

计算自己的加法份额：
\begin{equation}\label{eq:06-rvole:6}
y = -\sum_j g_j \cdot \alpha^0_j \pmod n.
\end{equation}

Sender 发送 $\tilde{a}_j$，本地保存 $y$。

\subsubsection{Receiver 获得加法份额}

Receiver 对每个 $j$ 计算自己的份额：
\begin{align}t_j &:= \alpha^{\beta_j}_j+\beta_j\cdot\tilde{a}_j \pmod n, \\
\textrm{a.k.a.~} t_j &= \alpha^0_j + \beta_j \cdot w \pmod n.
\end{align}

Receiver 计算 $z$ 份额：
\begin{equation}\label{eq:06-rvole:7}
z := \sum_j g_j \cdot t_j \pmod{n}.
\end{equation}

本章协议结束。

\begin{center}\rule{0.5\linewidth}{0.5pt}\end{center}

\section{RVOLE 协议完全版}

恶意 Sender 可能对不同的 OT 实例 $j$ 使用不同的 $w' \neq w$，
破坏 MtA 关系 $y + z = w \cdot \beta \pmod n$ 的正确性。
为此，我们约定每个 OT 实例携带 $N_2$ 个负载用于校验。
$N_2$ 对应原文的参数 $\rho$，取值由下述安全参数决定，并非固定为 1。
决定引入校验维度以后，将使 RVOLE 内部运算变得向量化。既然如此，
不妨继续引入 $(N_1-1)$ 个 MtA 关系。对于 ECDSA 签名，有 $N_1=2$。

综上，完全版的协议兑现如下 $N_1$ 个 MtA 关系：
\begin{equation}\label{eq:06-rvole:batch-mta}
y_{i}+z_{i} := w_{i}\cdot \beta \pmod n,
\quad i \in [0,\, N_1). \tag{batch-mta}
\end{equation}

\subsection{参数与编码}
\label{sec:rvole:parameters-model}

【参数】本文标量域为 $\mathbb Z_n$，其中 $n$ 是素数阶，$\kappa=\lceil\log_2 n\rceil$。
$\lambda_c$ 与 $\lambda_s$ 分别是计算安全参数和统计安全参数，不能与模数位宽混为一谈。
\href{https://eprint.iacr.org/2023/765.pdf}{DKLS23 Protocol 5.2} 的记号与本文对应如下：
\begin{center}
\begin{tabular}{lll}
\toprule
含义 & 原文 & 本文 \\
\midrule
标量模数 & $q$ & $n$ \\
功能负载数 & $\ell$ & $N_1$ \\
OT 实例数 & $\xi$ & $L=\kappa+2\lambda_s$ \\
校验负载数 & $\rho$ & $N_2=\lceil\kappa/\lambda_c\rceil$ \\
\bottomrule
\end{tabular}
\end{center}
本文以 $\kappa=256$、$\lambda_c=\lambda_s=128$ 为具体参数示例，得到 $L=512$、$N_2=2$。
ECDSA 同时计算两个乘积，因此 $N_1=2$，每个 OT 实例共携带四个域元素。
若改取 $\lambda_c=\kappa=256$，按上述原文公式才有 $N_2=1$，哈希输出长度也须随之调整。
这些是 RVOLE 的参数设置；提高其中一个参数并不会提高其他组件或整套 ECDSA 的安全强度。

【域与采样】所有 $w_k,\alpha^b_{j,k},\tilde a_{j,k},\theta_{k,i},\eta_k,\mu_{j,k},y_i,z_i$
均为 $\mathbb Z_n$ 元素，所有标量运算均模 $n$；最终摘要 $\mu$ 是比特串。
理想 OT 扩展接口直接提供独立均匀的域元素向量。
短 OT 密钥 $\rho_j^b$ 是实现该接口的种子，本文用以它为密钥的 $\mathrm{PRF}^{\mathbb Z_n}_{\rho_j^b}$ 表示域值扩展，
同一实例的 Receiver 只计算所选种子的扩展值。
实际扩展可由带域分离的密码学伪随机流实现：逐次读取 $\kappa$ 比特，解释为整数，拒绝不小于 $n$ 的值。
在理想随机流下，该拒绝采样精确均匀；实际密码学扩展只提供计算上的伪随机性。
因此，不能将「理想接口中的独立均匀」直接称为实际短种子扩展的无条件独立均匀。

【公开随机量】gadget 向量按 $\boldsymbol g\stackrel{\$}{\leftarrow}\mathbb Z_n^L$ 产生并公开，
采用原文 Protocol 5.2 脚注 a 的可信公共随机串设置，可跨实例复用。
若以公共种子和哈希实现，则须采用域分离和上述域采样，并在相应随机预言机模型下分析。
挑战 $\theta$ 是由公开 transcript 决定的随机预言机输出，双方都能重算，不是带秘密密钥的 PRF。
下文以 $\mathrm{RO}_{\mathbb Z_n^{N_2\times N_1}}$ 表示均匀域矩阵输出，
以 $\mathrm{RO}_{\{0,1\}^{2\lambda_c}}$ 表示用于压缩校验材料的摘要输出。
两种用途分别使用标签 \texttt{"challenge"} 与 \texttt{"check"}；示例参数下摘要长 256 比特。
参见\href{https://eprint.iacr.org/2023/765.pdf}{Protocol 5.2 第 3～4 步}。

【编码】$\mathrm{Enc}$ 表示无歧义编码，包含协议版本、域与维度参数，以及各字段的长度或类型。
域元素以 $[0,n)$ 的代表元作定长大端编码；数组按索引递增排列，矩阵先行后列。
具体地，$\tilde a$ 按 $j$、$k$ 排列，$\eta$ 按 $k$ 排列，校验矩阵按 $j$、$k$ 排列。
数组不得按无序集合序列化；消息中的维度、元素范围和编码不合法时即拒绝。
$\mathtt{sid}$ 唯一标识含有向双方身份的 RVOLE 实例，重试使用新实例号与新选择串。

\subsection{规格}

【输入】
Sender 持有 $N_1$ 个秘密因子 $w_i$，$i \in [0, N_1)$；
在协议里现场采集 $N_2$ 个新鲜随机校验因子 $w_k$，$k \in [N_1, N_1+N_2)$。
Receiver 持有秘密选项串 $\vec\beta \in \left\{0,1\right\}^L$，
对应着秘密因子 $\beta=\sum_{j\in[0,L)} \beta_j \cdot g_j$。
注意 $\vec\beta$ 的长度大于 $\beta$ 的比特数。

【输出】
Sender 获得 $N_1$ 个秘密加项 $y_i$。
Receiver 获得 $N_1$ 个秘密加项 $z_i$，使得等式 $\eqref{eq:06-rvole:batch-mta}$ 成立。

\subsection{实施}

协议有 $L$ 个 OT 实例，用 $j$ 索引。每个 OT 实例携带 $N_1+N_2$ 个负载，用 $k$ 或 $i$ 索引。

\subsubsection{Sender 获得加法份额}

对每个 OT 实例 $j\in [0, L)$ 和负载 $k\in[0, N_1+N_2)$，
Sender 首先对 OT 密钥 $\rho_j^b$ 进行衍生，
得到成对拓展密钥 如公式 \eqref{eq:06-rvole:extkey}。
然后计算密钥公差用 $w_k$ 加密后的密文，如公式 \eqref{eq:06-rvole:offset}。
最后对同样的 $j$ 和每个功能负载 $i \in [0, N_1)$，
Sender 计算加法份额，如公式 \eqref{eq:06-rvole:yi}。
\begin{align}\label{eq:06-rvole:extkey}
\tag{extkey}
\alpha^b_{j,k} &:= \mathrm{PRF}^{\mathbb Z_n}_{\rho_j^b}\bigl(
\mathrm{Enc}(\mathtt{sid},\texttt{"ot-pad"},j,k)
\bigr)\in\mathbb Z_n,
b\in\left\{0,1\right\}.
\\
\tilde{a}_{j,k} &:= \alpha^0_{j,k}-\alpha^1_{j,k}+w_{k}.
\\
\label{eq:06-rvole:offset}
\tag{offset}
\tilde a &:= (\tilde a_{j,k})_{j\in[0,L),\,k\in[0,N_1+N_2)}.
\\
\label{eq:06-rvole:yi}
\tag{yi}
y_{i} &:= -\sum_j g_j\cdot \alpha^0_{j,i} ~.
\end{align}

至此 Sender 已得到全部的输出 $y_i, i\in[0, N_1)$。

\subsubsection{Sender 响应挑战，发送消息}
\label{sec:rvole:challenge-response}

在本节剩余内容里，$j\in[0, L)$，也就是索引每个 OT实例；
$k\in[N_1, N_1+N_2)$，也就是索引每个校验负载。

对每个 $k$，
Sender 首先从已知信息派生出挑战向量 $\theta_k$，如公式 \eqref{eq:06-rvole:ch}。
然后计算线性响应，如公式 \eqref{eq:06-rvole:resp-lin}，
以及哈希响应，如公式 \eqref{eq:06-rvole:resp-ha}。
\begin{align}
\label{eq:06-rvole:ch}\tag{ch}
(\theta_k)_{k\in[N_1,N_1+N_2)}
&:= \mathrm{RO}_{\mathbb Z_n^{N_2\times N_1}}\bigl(
\mathrm{Enc}(\mathtt{sid},\texttt{"challenge"},\tilde a)\bigr).
\\
\label{eq:06-rvole:resp-lin}\tag{resp-lin}
\eta_k &:= w_{k}+\sum_{i\in[0,N_1)}\theta_{k,i}\cdot w_{i} \pmod n,
\\
\eta &:= (\eta_k)_{k\in[N_1,N_1+N_2)}.
\\
\mu_{j,k} &:= \alpha^0_{j,k}
  +\sum_{i\in[0,N_1)}\theta_{k,i}\cdot \alpha^0_{j,i},
\\
\label{eq:06-rvole:resp-ha}\tag{resp-ha}
\mu &:= \mathrm{RO}_{\{0,1\}^{2\lambda_c}}\bigl(
\mathrm{Enc}(\mathtt{sid},\texttt{"check"},(\mu_{j,k})_{j,k})\bigr).
\end{align}

最后，Sender 发送 $\tilde a$，$\eta$，$\mu$。

\subsubsection{Receiver 验证挑战，获得加法份额}

在本节，$j\in[0, L)$ 索引每个 OT 实例；$i\in[0, N_1)$ 索引每个功能负载。

对每个 $j$ 和全部负载 $k \in [0, N_1+N_2)$，
Receiver 对 OT 密钥 $\rho_j^{\beta_j}$ 进行如公式 \eqref{eq:06-rvole:extkey} 
的衍生（$b=\beta_j$），然后用收到的 $\tilde a_{j,k}$ 计算
\begin{align}
t_{j,k} &:= \alpha^{\beta_j}_{j,k}+\beta_j\cdot \tilde{a}_{j,k}, \\
\tag{tjk}\label{eq:06-rvole:tjk}
&= \alpha^0_{j,k}+\beta_j\cdot w_{k}~,
\end{align}

先按公式 \eqref{eq:06-rvole:ch} 重放完整的挑战矩阵 $(\theta_k)_k$，所有 OT 实例共享该矩阵。
之后，对每个 $j$ 和每个校验负载 $k$，计算公式 \eqref{eq:06-rvole:mu-hat}，
检验 $\hat\mu\stackrel{?}{=}\mu$。
如果检验通过，那么 Receiver 对每个功能负载 $i$ 计算本地份额如公式 \eqref{eq:06-rvole:zi}。
\begin{align}
\hat\mu_{j,k} &:= t_{j,k}
+ \sum_{i\in[0,N_1)} \theta_{k,i}\cdot t_{j,i}
- \beta_j\cdot\eta_k, 
\\
\label{eq:06-rvole:mu-hat}
\hat\mu &:= \mathrm{RO}_{\{0,1\}^{2\lambda_c}}\bigl(
  \mathrm{Enc}(\mathtt{sid},\texttt{"check"},(\hat\mu_{j,k})_{j,k})
\bigr).
\\
\label{eq:06-rvole:zi}\tag{zi}
z_{i} &:= \sum_{j}g_j\cdot t_{j,i} ~.
\end{align}

最后，Receiver 本地保存 $z_{i}$。

\begin{center}\rule{0.5\linewidth}{0.5pt}\end{center}

\section{这种挑战是怎么想到的?}

本附录复盘 \hyperref[sec:rvole:challenge-response]{“挑战-响应”}的设计过程，展示它不是天降神来之笔、注意力惊人的奇迹，
而是由若干“思维脚手架”搭建出来的。

\subsection{给出作恶/诚实的形式化描述}

Sender 在本协议只发一条消息 $(\tilde a, \eta, \mu)$。
其中 $\eta, \mu$ 恰是被检查的对象，所以真正的作弊自由度全在 $\tilde a$ 里。
由公式 \eqref{eq:06-rvole:offset}，Sender 知道全部 $\alpha^0_{j,k}, \alpha^1_{j,k}$，
所以他发出的任意 $\tilde a_{j,k}$ 都可以反解出一个“等效输入”如下式
\begin{equation}\label{eq:06-rvole:16}
w_{j,k} := \tilde a_{j,k} - \alpha^0_{j,k} + \alpha^1_{j,k} \pmod n.
\end{equation}

至此，“Sender 是诚实的”，当且仅当
“对每个 $k$，都有 $\forall (j_1, j_2), w_{j_1,k}=w_{j_2,k}$”。

\subsection{用“公共基准”回答“两两相等”}

Receiver 能直接比较两行吗？他只能把两行相减，如公式 \eqref{eq:06-rvole:diff}。
式中第二个括号项是消不掉的，因而直接对比的思路行不通。
\begin{equation}
\label{eq:06-rvole:diff}
t_{j_1,k} - t_{j_2,k} = \left(\alpha^0_{j_1,k} - \alpha^0_{j_2,k}\right) + \left(\beta_{j_1} w_{j_1,k} - \beta_{j_2} w_{j_2,k}\right),
\end{equation}

两行的 pad 是各自独立的随机数，消不掉。直接对比的思路行不通。问题化归：
一列内部两两相等 iff 全列等于同一个公共值。

问题化归：“一列内部两两相等” 当且仅当 “一列等于同一个公共值 $c_k$”。
化归以后，需要让 Sender 申报这个公共值 $c_k$。
若第 $j$ 行确实用了申报值 $c_k$，则根据公式 \eqref{eq:06-rvole:tjk} 有
\begin{equation}
t_{j,k} - \beta_j \cdot c_k = \alpha^0_{j,k}.
\end{equation}
若 Sender 实际用的是 $c_k + e$（$e \neq 0$），则 Receiver 的实算值 = 预测值 $+ \beta_j e$。
此时 Sender 要想蒙混过关，就必须猜 $\beta_j$。
这就是整个检查的骨架：诚实时两侧恒等且与 $\vec\beta$ 无关，
作弊时差一个 $\beta_j$ 相关项，Sender 猜不中
（类似于 KOS15 抓恶意 Receiver 的思路）。

到此，“申报 + 逐行对照”已经能检查一列。
剩两个问题：① 功能列的公共值就是秘密 $w_i$，不能明着申报；② 申报必须先于挑战固定。

\subsection{随机线性组合}

$w_i, i\in[0,N_1)$ 不能明文申报。一个经典技巧是申报所有 $w_i$ 的随机线性组合，
这恰好就是 $\eta_k$，如公式 \eqref{eq:06-rvole:resp-lin}。
该式的关键性质：
对随机 $\theta$，“组合列是常数列” 以压倒性概率当且仅当 “每个参与列都是常数列”。

\subsection{牺牲负载}

公式 \eqref{eq:06-rvole:resp-lin} 为什么要混入一个加项 $w_k$？
若组合只取功能列，则 Sender 申报的是 $\eta_k = \sum_i \theta_{k,i} w_i$；
考虑到 $\theta$ 是公开的系数，这就泄露关于所有 $w_i$ 的一个线性方程。
标准补丁是牺牲（sacrifice）：加一个新鲜随机的 $w_k$ 进组合，充当一次一密掩码。

这个技巧的妙处在于，Receiver 手里有 $t_{j,k} = \alpha^0_{j,k} + \beta_j w_k$，
这意味着 Receiver 能在本地消掉 $\beta_j w_k$。

\subsection{Fiat-Shamir 定序}

$\theta$ 必须在 Sender 固定 $\tilde a$（等价于固定全部 $w_{j,k}$）
之后才可知，否则 Sender 可以挑选对他有利的 $\theta$。
采用 Fiat-Shamir 之后，作弊 Sender 想过检只有两条路：

\begin{itemize}
  \item \textbf{研磨（grinding）}：
  这是一条跟哈希安全性死磕的路。
  反复重选 $\tilde a$，直到摇出的 $\theta$ 恰好使所有偏差组合为零。
  \item \textbf{猜比特（selective failure）}：
  这是一条绕过哈希安全性的路。
  Sender 不再期待猜中整个 $\vec\beta$，只期待猜中某些 $\beta_j$。
  每一位猜中的概率为 $1/2$，猜错即 abort。
  注意到 Receiver 的实际秘密为标量 $\beta=\sum_{j=0}^{L-1} \beta_j \cdot g_j$，
  而参与 OT 的是更长的比特向量 $\vec\beta=(\beta_0,\ldots,\beta_{L-1})$；
  这意味着 gadget 稀释了 $\beta$ 的信息量，从而提高 selective failure 的难度。
\end{itemize}
